\oldpage[I]{89}
ainsi obtenus vérifient les conditions de compatibilité avec les opérateurs de
restriction (\hyperref[0.1.4.2-fr]{0, 1.4.2}), donc définissent un isomorphisme
canonique fonctoriel
\[
  \widetilde{M}\otimes_{\widetilde{A}}\widetilde{N}
  \simeq\widetilde{(M\otimes_A N)}
\]

(ii) Le faisceau
$\mathscr{G}=\mathscr{H}\!om_{\widetilde{A}}(\widetilde{M},\widetilde{N})$
est associé au préfaisceau
\[
  U\longmapsto\mathscr{G}(U)
  =\operatorname{Hom}_{\widetilde{A}|U}
  (\widetilde{M}|U,\widetilde{N}|U)
\]
$U$ parcourant la base de $X$ formée des $D(f)$. Or,
$\mathscr{G}(D(f))$ s'identifie canoniquement à
$\operatorname{Hom}_{A_f}(M_f,N_f)$
(\hyperref[I.1.3.6-fr]{1.3.6} et \hyperref[I.1.3.8-fr]{1.3.8}), qui
lui-même, en vertu de l'hypothèse sur $M$, s'identifie canoniquement à
$(\operatorname{Hom}_A(M,N))_f$
(\hyperref[0.1.3.5-fr]{0, 1.3.5}). Finalement,
$(\operatorname{Hom}_A(M,N))_f$ s'identifie canoniquement à
$\Gamma(D(f),\widetilde{(\operatorname{Hom}_A(M,N))})$
(\hyperref[I.1.3.6-fr]{1.3.6} et \hyperref[I.1.3.7-fr]{1.3.7}), et les
isomorphismes canoniques
$\mathscr{G}(D(f))\simeq
\Gamma(D(f),\widetilde{(\operatorname{Hom}_A(M,N))})$ ainsi obtenus sont
compatibles avec les opérateurs de restriction
(\hyperref[0.1.4.2-fr]{0, 1.4.2})~; ils définissent donc un isomorphisme
canonique
\[
  \mathscr{H}\!om_{\widetilde{A}}(\widetilde{M},\widetilde{N})
  \simeq\widetilde{(\operatorname{Hom}_A(M,N))}.
\]

\begin{env}[1.3.13]
\label{I.1.3.13-fr}
Soit maintenant $B$ une $A$-algèbre (commutative)~; cela peut s'interpréter
en disant que $B$ est un $A$-module et qu'on s'est donné un élément $e\in B$
et un $A$-homomorphisme
$\varphi:B\otimes_A B\to B$, de sorte que~: $1^\circ$ Les diagrammes
\[
  \xymatrix{
    B\otimes_A B\otimes_A B\ar[r]^{\varphi\otimes 1}
      \ar[d]_{1\otimes\varphi} &
    B\otimes_A B\ar[d]^\varphi & &
    B\otimes_A B\ar[rr]^\sigma\ar[rd]_\varphi & &
    B\otimes_A B\ar[dl]^\varphi\\
    B\otimes_A B\ar[r]_\varphi &
    B & & & B
  }
\]
($\sigma$ symétrie canonique) soient commutatifs~; $2^\circ$
$\varphi(e\otimes x)=\varphi(x\otimes e)=x$. En vertu de
(\hyperref[I.1.3.12-fr]{1.3.12}), l'homomorphisme
$\widetilde{\varphi}:\widetilde{B}\otimes_{\widetilde{A}}\widetilde{B}
\to\widetilde{B}$ de $\widetilde{A}$-Modules vérifie les conditions
analogues, donc définit sur $\widetilde{B}$ une structure de
$\widetilde{A}$-Algèbre. De la même manière, la donnée d'un $B$-module $N$
revient à la donnée d'un $A$-module $N$ et d'un $A$-homomorphisme
$\psi:B\otimes_A N\to N$ tel que le diagramme
\[
  \xymatrix{
    B\otimes_A B\otimes_A N\ar[r]^{\varphi\otimes 1}
      \ar[d]_{1\otimes\psi} &
    B\otimes_A N\ar[d]^\psi\\
    B\otimes_A N\ar[r]_\psi & N
  }
\]
soit commutatif et $\psi(e\otimes n)=n$~; l'homomorphisme
$\widetilde{\psi}:\widetilde{B}\otimes_{\widetilde{A}}\widetilde{N}
\to\widetilde{N}$ vérifiera la condition analogue, et définit donc sur
$\widetilde{N}$ une structure de $\widetilde{B}$-Module.

De la même façon, on voit que si $u:B\to B'$ (resp. $v:N\to N'$) est un
homomorphisme de $A$-algèbres (resp. de $B$-modules), $\widetilde{u}$
(resp. $\widetilde{v}$) est un homomorphisme de $\widetilde{A}$-Algèbres
(resp. de $\widetilde{B}$-Modules), $\operatorname{Ker}\widetilde{u}$ un
$\widetilde{B}$-Idéal (resp. $\operatorname{Ker}\widetilde{v}$,
$\operatorname{Coker}\widetilde{v}$ et $\operatorname{Im}\widetilde{v}$
des $\widetilde{B}$-Modules). Si $N$ est un $B$-module, $\widetilde{N}$ est
un $\widetilde{B}$-Module de type fini si et seulement si $N$ est un
$B$-module de type fini (\hyperref[0.5.2.3-fr]{0, 5.2.3}).
